%%%%%%%%%%%%%%%%%%%%%%%%%%%%%%%%%%%%%%%%%%%%%%%%%%%%%%%%%%%%%%%%%%%
% Header: General study / work / project template
%%%%%%%%%%%%%%%%%%%%%%%%%%%%%%%%%%%%%%%%%%%%%%%%%%%%%%%%%%%%%%%%%%%

\documentclass[11pt,a4paper]{scrartcl}

% Layout
\usepackage[margin=2.5cm]{geometry}
\usepackage{setspace}
\setstretch{1.15}

% General packages
\usepackage{graphicx}
\usepackage{enumitem}
\usepackage[most]{tcolorbox}
\tcbuselibrary{breakable}
\usepackage[breaklinks=true,colorlinks=true,linkcolor=blue,urlcolor=blue,citecolor=blue]{hyperref}
\usepackage{amsmath,amsthm,amssymb,mathtools}
\usepackage{braket}
\usepackage{icomma}
\usepackage[T1]{fontenc}
\usepackage[utf8]{inputenc}
\usepackage[ngerman]{babel}

% Units / tables / floats / references
\usepackage[locale = DE,space-before-unit=true,per-mode = symbol]{siunitx}
\usepackage{booktabs,multirow}
\usepackage{placeins}
\usepackage[natbib,abbreviate=true,doi=false,style=numeric-comp,giveninits=true,sorting=none]{biblatex}
\usepackage{csquotes}

% Optional math shortcuts
\newcommand{\R}{\mathbb{R}}
\newcommand{\N}{\mathbb{N}}
\newcommand{\Q}{\mathbb{Q}}
\newcommand{\abs}[1]{\left| #1 \right|}
\newcommand{\norm}[1]{\left\lVert #1 \right\rVert}
\newcommand{\ol}[1]{\overline{#1}}
\newcommand{\op}[1]{\hat{#1}}
\newcommand{\Eone}{E_1^{(0)}}
\newcommand{\Etwo}{E_2^{(0)}}
\newcommand{\DeltaE}{\Delta}

% Box styles
\newtcolorbox{calculationbox}[1][]{
    title=Calculation,
    breakable,
    enhanced,
    colback=white,
    colframe=black!65,
    #1
}

\newtcolorbox{solutionbox}[1][]{
    title=Solution,
    breakable,
    enhanced,
    colback=white,
    colframe=black!65,
    #1
}

\newtcolorbox{answerbox}[1][]{
    title=Answer,
    breakable,
    enhanced,
    colback=white,
    colframe=black!65,
    #1
}

\newtcolorbox{notebox}[1][]{
    title=Notes,
    breakable,
    enhanced,
    colback=white,
    colframe=black!65,
    #1
}

\newtcolorbox{sidecalcbox}[1][]{
    title=Side Calculation,
    breakable,
    enhanced,
    colback=white,
    colframe=black!65,
    #1
}

%%%%%%%%%%%%%%%%%%%%%%%%%%%%%%%%%%%%%%%%%%%%%%%%%%%%%%%%%%%%%%%%%%%
\begin{document}

    \title{Theoretische Physik 2 -- Blatt 09}
    \author{Tobias Laser}
    \date{\today}
    \maketitle
    \vfill
    \thispagestyle{empty}
    \newpage

    \pagenumbering{arabic}

    % ================================================================
    \paragraph{21. Zeitunabhängige Störungstheorie und ein 2-Niveau-System:}

    Wir betrachten ein 2-Niveau-System mit Hamiltonoperator
    \begin{align*}
        \op H
        =
        \op H_0 + \lambda \op V
        =
        \begin{pmatrix}
            E_1^{(0)} & 0 \\
            0 & E_2^{(0)}
        \end{pmatrix}
        +
        \lambda
        \begin{pmatrix}
            1 & 1 \\
            1 & 1
        \end{pmatrix}.
    \end{align*}

    \begin{enumerate}[label=(\alph*)]
        \item Finden Sie die Eigenzustände von $\op H_0$ (ohne Rechnung), sowie eine exakte Lösung für die Eigenwerte $E_1,E_2$ und Eigenzustände $\ket{\psi_1},\ket{\psi_2}$ des gesamten Hamiltonoperators $\op H$.

        \begin{calculationbox}
            Da $\op H_0$ bereits diagonal ist, sind die Eigenzustände von $\op H_0$ direkt die beiden Standardbasiszustände. Für den gesamten Hamiltonoperator addiert die Störung auf jedes Matrixelement den Beitrag $\lambda$ dort, wo in $V$ eine $1$ steht. Deshalb lautet die Matrix von $\op H$
            \begin{align*}
                \op H
                =
                \begin{pmatrix}
                    E_1^{(0)}+\lambda & \lambda\\
                    \lambda & E_2^{(0)}+\lambda
                \end{pmatrix}.
            \end{align*}

            \begin{solutionbox}[title=Ungestörte Eigenzustände]
                \begin{align*}
                    \ket{1}=\begin{pmatrix}1\\0\end{pmatrix},
                    \qquad
                    \ket{2}=\begin{pmatrix}0\\1\end{pmatrix},
                    \qquad
                    \op H_0\ket{1}=E_1^{(0)}\ket{1},
                    \qquad
                    \op H_0\ket{2}=E_2^{(0)}\ket{2}.
                \end{align*}
            \end{solutionbox}

            \begin{calculationbox}[title=Exakte Eigenwerte]
                Für die exakten Eigenwerte lösen wir die charakteristische Gleichung. Dazu setzen wir
                \begin{align*}
                    \det(\op H-E\mathbb 1)=0.
                \end{align*}
                Damit folgt
                \begin{align*}
                    0
                    &=
                    \det
                    \begin{pmatrix}
                        E_1^{(0)}+\lambda-E & \lambda\\
                        \lambda & E_2^{(0)}+\lambda-E
                    \end{pmatrix}\\
                    &=
                    (E_1^{(0)}+\lambda-E)(E_2^{(0)}+\lambda-E)-\lambda^2.
                \end{align*}
                Um die quadratische Gleichung übersichtlich zu schreiben, verwenden wir
                \begin{align*}
                    \bar E^{(0)}=\frac{E_1^{(0)}+E_2^{(0)}}{2},
                    \qquad
                    \Delta=E_1^{(0)}-E_2^{(0)}.
                \end{align*}

                \begin{solutionbox}[title=Exakte Eigenwerte]
                    \begin{align*}
                        E_{\pm}
                        =
                        \bar E^{(0)}+\lambda
                        \pm
                        \sqrt{\left(\frac{\Delta}{2}\right)^2+\lambda^2}.
                    \end{align*}
                \end{solutionbox}
            \end{calculationbox}

            \begin{calculationbox}[title=Exakte Eigenzustände]
                Wir schreiben einen Eigenzustand als Linearkombination der Basiszustände:
                \begin{align*}
                    \ket{\psi_\pm}=a_\pm\ket{1}+b_\pm\ket{2}.
                \end{align*}
                Die erste Zeile der Eigenwertgleichung $(\op H-E_\pm)\ket{\psi_\pm}=0$ liefert
                \begin{align*}
                    (E_1^{(0)}+\lambda-E_\pm)a_\pm+\lambda b_\pm=0.
                \end{align*}
                Daher kann man
                \begin{align*}
                    b_\pm=\frac{E_\pm-E_1^{(0)}-\lambda}{\lambda}a_\pm
                \end{align*}
                wählen. Wir definieren den Quotienten
                \begin{align*}
                    q_\pm
                    :=
                    \frac{b_\pm}{a_\pm}
                    =
                    \frac{E_\pm-E_1^{(0)}-\lambda}{\lambda}
                    =
                    \frac{-\Delta/2\pm\sqrt{(\Delta/2)^2+\lambda^2}}{\lambda}.
                \end{align*}

                \begin{solutionbox}[title=Exakte Eigenzustände]
                    \begin{align*}
                        \ket{\psi_\pm}
                        =
                        \frac{1}{\sqrt{1+q_\pm^2}}
                        \left(\ket{1}+q_\pm\ket{2}\right),
                        \qquad
                        q_\pm=
                        \frac{-\Delta/2\pm\sqrt{(\Delta/2)^2+\lambda^2}}{\lambda}.
                    \end{align*}
                \end{solutionbox}
            \end{calculationbox}
        \end{calculationbox}

        \item Nehmen Sie an, dass $\lambda \ll \abs{E_1^{(0)}-E_2^{(0)}}$, und berechnen Sie mit zeitunabhängiger Störungstheorie die Energieeigenwerte in erster und zweiter Ordnung in $\lambda$, sowie die Energieeigenzustände bis zu erster Ordnung in $\lambda$. Vergleichen Sie dieses Ergebnis mit (a).

        \begin{calculationbox}
            Hier ist der Abstand der ungestörten Energien groß gegenüber der Störung. Deshalb ist nicht-entartete Störungstheorie anwendbar. Wir setzen wieder
            \begin{align*}
                \Delta=E_1^{(0)}-E_2^{(0)}.
            \end{align*}
            Die relevanten Matrixelemente der Störung sind
            \begin{align*}
                V_{11}=V_{22}=V_{12}=V_{21}=1.
            \end{align*}

            \begin{calculationbox}[title=Energien bis zweiter Ordnung]
                Für einen nicht-entarteten Zustand $\ket n$ gilt
                \begin{align*}
                    E_n
                    =
                    E_n^{(0)}
                    +\lambda V_{nn}
                    +\lambda^2\sum_{m\neq n}
                    \frac{\abs{V_{mn}}^2}{E_n^{(0)}-E_m^{(0)}}
                    +\mathcal O(\lambda^3).
                \end{align*}
                Für $\ket 1$ erhalten wir
                \begin{align*}
                    E_1
                    &=
                    E_1^{(0)}
                    +\lambda V_{11}
                    +\lambda^2\frac{\abs{V_{21}}^2}{E_1^{(0)}-E_2^{(0)}}
                    +\mathcal O(\lambda^3)\\
                    &=
                    E_1^{(0)}
                    +\lambda
                    +\frac{\lambda^2}{\Delta}
                    +\mathcal O(\lambda^3).
                \end{align*}
                Für $\ket 2$ erhalten wir analog
                \begin{align*}
                    E_2
                    &=
                    E_2^{(0)}
                    +\lambda V_{22}
                    +\lambda^2\frac{\abs{V_{12}}^2}{E_2^{(0)}-E_1^{(0)}}
                    +\mathcal O(\lambda^3)\\
                    &=
                    E_2^{(0)}
                    +\lambda
                    -\frac{\lambda^2}{\Delta}
                    +\mathcal O(\lambda^3).
                \end{align*}

                \begin{solutionbox}[title=Energieeigenwerte für $\lambda\ll\abs\Delta$]
                    \begin{align*}
                        E_1
                        =
                        E_1^{(0)}+\lambda+\frac{\lambda^2}{\Delta}
                        +\mathcal O(\lambda^3),
                        \qquad
                        E_2
                        =
                        E_2^{(0)}+\lambda-\frac{\lambda^2}{\Delta}
                        +\mathcal O(\lambda^3).
                    \end{align*}
                \end{solutionbox}
            \end{calculationbox}

            \begin{calculationbox}[title=Eigenzustände bis erster Ordnung]
                Für die Korrektur eines nicht-entarteten Eigenzustands gilt
                \begin{align*}
                    \ket{n^{(1)}}
                    =
                    \lambda
                    \sum_{m\neq n}
                    \frac{V_{mn}}{E_n^{(0)}-E_m^{(0)}}\ket m.
                \end{align*}
                Also folgt für $\ket1$
                \begin{align*}
                    \ket{\psi_1}
                    =
                    \ket1+\lambda\frac{V_{21}}{E_1^{(0)}-E_2^{(0)}}\ket2
                    +\mathcal O(\lambda^2)
                    =
                    \ket1+\frac{\lambda}{\Delta}\ket2+\mathcal O(\lambda^2).
                \end{align*}
                Für $\ket2$ folgt
                \begin{align*}
                    \ket{\psi_2}
                    =
                    \ket2+\lambda\frac{V_{12}}{E_2^{(0)}-E_1^{(0)}}\ket1
                    +\mathcal O(\lambda^2)
                    =
                    \ket2-\frac{\lambda}{\Delta}\ket1+\mathcal O(\lambda^2).
                \end{align*}

                \begin{solutionbox}[title=Eigenzustände bis erster Ordnung]
                    \begin{align*}
                        \ket{\psi_1}
                        =
                        \ket1+\frac{\lambda}{\Delta}\ket2+\mathcal O(\lambda^2),
                        \qquad
                        \ket{\psi_2}
                        =
                        \ket2-\frac{\lambda}{\Delta}\ket1+\mathcal O(\lambda^2).
                    \end{align*}
                \end{solutionbox}
            \end{calculationbox}

            \begin{calculationbox}[title=Vergleich mit der exakten Lösung]
                Für $\abs\lambda\ll\abs\Delta$ wird die Wurzel aus Teil (a) entwickelt. Für $\Delta>0$ gilt
                \begin{align*}
                    \sqrt{\left(\frac{\Delta}{2}\right)^2+\lambda^2}
                    &=
                    \frac{\Delta}{2}\sqrt{1+\frac{4\lambda^2}{\Delta^2}}\\
                    &=
                    \frac{\Delta}{2}
                    \left(1+\frac{2\lambda^2}{\Delta^2}+\mathcal O(\lambda^4)\right)\\
                    &=
                    \frac{\Delta}{2}+\frac{\lambda^2}{\Delta}+\mathcal O(\lambda^4).
                \end{align*}
                Einsetzen in $E_\pm$ liefert genau die Energien aus der Störungstheorie. Die Entwicklung der exakten Eigenzustände liefert ebenfalls dieselben ersten Ordnungen.

                \begin{solutionbox}[title=Vergleich]
                    \begin{align*}
                        E_+=E_1^{(0)}+\lambda+\frac{\lambda^2}{\Delta}+\mathcal O(\lambda^4),
                        \qquad
                        E_-=E_2^{(0)}+\lambda-\frac{\lambda^2}{\Delta}+\mathcal O(\lambda^4).
                    \end{align*}
                \end{solutionbox}
            \end{calculationbox}
        \end{calculationbox}

        \item Nehmen Sie an, dass $E_1^{(0)}$ und $E_2^{(0)}$ fast entartet sind, d.h. $\abs{E_1^{(0)}-E_2^{(0)}}\ll \lambda$, und berechnen Sie die Energieeigenwerte bis zu erster Ordnung in $\lambda$, sowie die Energieeigenzustände nullter Ordnung mit Störungstheorie im Fall $E_1^{(0)}=E_2^{(0)}$. Vergleichen Sie das Ergebnis mit der exakten Lösung von (a).

        \begin{calculationbox}
            Wenn die ungestörten Energien entartet oder fast entartet sind, ist die nicht-entartete Formel nicht geeignet, weil die Nenner klein werden. Im exakt entarteten Fall setzen wir
            \begin{align*}
                E_1^{(0)}=E_2^{(0)}=:E^{(0)}.
            \end{align*}
            Dann muss die Störung im entarteten Unterraum diagonalisiert werden:
            \begin{align*}
                V=
                \begin{pmatrix}
                    1&1\\
                    1&1
                \end{pmatrix}.
            \end{align*}

            \begin{calculationbox}[title=Diagonalisierung im entarteten Unterraum]
                Für den symmetrischen Zustand gilt
                \begin{align*}
                    V\frac{1}{\sqrt2}\begin{pmatrix}1\\1\end{pmatrix}
                    =
                    \frac{1}{\sqrt2}\begin{pmatrix}2\\2\end{pmatrix}
                    =
                    2\frac{1}{\sqrt2}\begin{pmatrix}1\\1\end{pmatrix}.
                \end{align*}
                Für den antisymmetrischen Zustand gilt
                \begin{align*}
                    V\frac{1}{\sqrt2}\begin{pmatrix}1\\-1\end{pmatrix}
                    =
                    \frac{1}{\sqrt2}\begin{pmatrix}0\\0\end{pmatrix}.
                \end{align*}

                \begin{solutionbox}[title=Richtige Zustände nullter Ordnung]
                    \begin{align*}
                        \ket{\psi_+^{(0)}}=\frac{1}{\sqrt2}(\ket1+\ket2),
                        \qquad
                        \ket{\psi_-^{(0)}}=\frac{1}{\sqrt2}(\ket1-\ket2).
                    \end{align*}
                \end{solutionbox}
            \end{calculationbox}

            \begin{calculationbox}[title=Energien erster Ordnung]
                Die Energieverschiebungen erster Ordnung sind die Eigenwerte von $\lambda V$. Also erhält der symmetrische Zustand die Verschiebung $2\lambda$ und der antisymmetrische Zustand die Verschiebung $0$.

                \begin{solutionbox}[title=Energien im entarteten Fall]
                    \begin{align*}
                        E_+=E^{(0)}+2\lambda,
                        \qquad
                        E_-=E^{(0)}.
                    \end{align*}
                \end{solutionbox}
            \end{calculationbox}

            \begin{calculationbox}[title=Vergleich mit der exakten Lösung]
                In der exakten Lösung aus Teil (a) setzen wir $\Delta=0$. Dann ist
                \begin{align*}
                    E_\pm=E^{(0)}+\lambda\pm\abs\lambda.
                \end{align*}
                Für $\lambda>0$ ergibt das
                \begin{align*}
                    E_+=E^{(0)}+2\lambda,
                    \qquad
                    E_-=E^{(0)}.
                \end{align*}
                Das stimmt mit der entarteten Störungstheorie überein. Der Unterschied zur nicht-entarteten Rechnung ist, dass man bei Entartung zuerst die Störung im entarteten Unterraum diagonalisiert.

                \begin{solutionbox}[title=Vergleich]
                    \begin{align*}
                        \Delta=0
                        \quad\Longrightarrow\quad
                        E_\pm=E^{(0)}+\lambda\pm\abs\lambda.
                    \end{align*}
                \end{solutionbox}
            \end{calculationbox}
        \end{calculationbox}
    \end{enumerate}

    % ================================================================
    \paragraph{22. Magnetische Dipolwechselwirkung:}

    Wir betrachten zwei Spin-$\frac12$-Teilchen mit magnetischem Moment
    \begin{align*}
        \vec M^{(i)}=\gamma\vec S^{(i)},
        \qquad i=1,2,
    \end{align*}
    in einem externen Magnetfeld
    \begin{align*}
        \vec B=B_0\vec e_z.
    \end{align*}
    Der Hamiltonoperator des Systems ist
    \begin{align*}
        \op H_0=-\vec B\cdot(\vec M^{(1)}+\vec M^{(2)}).
    \end{align*}

    \begin{enumerate}[label=(\alph*)]
        \item Zeigen Sie, dass der Hamiltonoperator $\op H_0$ diagonal in der Produktbasis
        \begin{align*}
            \ket{++},\quad \ket{+-},\quad \ket{-+},\quad \ket{--}
        \end{align*}
        ist und geben Sie die zugehörigen Energieeigenwerte an. Gibt es entartete Energieniveaus? Es ist nützlich, die Frequenz $\omega_0=\gamma B_0$ zu definieren.

        \begin{calculationbox}
            Da das Magnetfeld in $z$-Richtung zeigt, enthält $\op H_0$ nur $z$-Komponenten der Spins. Mit $\omega_0=\gamma B_0$ gilt
            \begin{align*}
                \op H_0
                =
                -B_0\gamma(\op S_z^{(1)}+\op S_z^{(2)})
                =
                -\omega_0(\op S_z^{(1)}+\op S_z^{(2)}).
            \end{align*}
            Die Produktbasis besteht aus Eigenzuständen von $\op S_z^{(1)}$ und $\op S_z^{(2)}$. Deshalb ist $\op H_0$ in dieser Basis diagonal.

            \begin{calculationbox}[title=Energieeigenwerte]
                Für Spin-$\frac12$ gilt
                \begin{align*}
                    \op S_z\ket+=\frac{\hbar}{2}\ket+,
                    \qquad
                    \op S_z\ket-=-\frac{\hbar}{2}\ket-.
                \end{align*}
                Daher
                \begin{align*}
                    \op H_0\ket{++}
                    &=-\omega_0\left(\frac{\hbar}{2}+\frac{\hbar}{2}\right)\ket{++}
                    =-\hbar\omega_0\ket{++},\\
                    \op H_0\ket{+-}
                    &=-\omega_0\left(\frac{\hbar}{2}-\frac{\hbar}{2}\right)\ket{+-}
                    =0,\\
                    \op H_0\ket{-+}
                    &=-\omega_0\left(-\frac{\hbar}{2}+\frac{\hbar}{2}\right)\ket{-+}
                    =0,\\
                    \op H_0\ket{--}
                    &=-\omega_0\left(-\frac{\hbar}{2}-\frac{\hbar}{2}\right)\ket{--}
                    =+\hbar\omega_0\ket{--}.
                \end{align*}

                \begin{solutionbox}[title=Ungestörtes Spektrum]
                    \begin{align*}
                        E_{++}^{(0)}=-\hbar\omega_0,
                        \qquad
                        E_{+-}^{(0)}=0,
                        \qquad
                        E_{-+}^{(0)}=0,
                        \qquad
                        E_{--}^{(0)}=+\hbar\omega_0.
                    \end{align*}
                \end{solutionbox}
                Das Niveau $E=0$ ist zweifach entartet, weil die Zustände $\ket{+-}$ und $\ket{-+}$ dieselbe Energie besitzen.
            \end{calculationbox}
        \end{calculationbox}

        \item Die Teilchen sind sich so nahe, dass die Dipol-Dipol-Wechselwirkung der beiden magnetischen Momente nicht vernachlässigt werden kann. Diese wird durch einen Hamiltonoperator
        \begin{align*}
            \op H_1
            =
            \frac{\mu_0}{4\pi}\frac{1}{r^3}
            \left[
                \vec M^{(1)}\cdot \vec M^{(2)}
                -3(\vec M^{(1)}\cdot\vec n)(\vec M^{(2)}\cdot\vec n)
            \right]
        \end{align*}
        beschrieben, wobei $r$ den Abstand zwischen den Teilchen und
        \begin{align*}
            \vec n=(\sin\theta\cos\varphi,\sin\theta\sin\varphi,\cos\theta)
        \end{align*}
        den Einheitsvektor in diese Richtung bezeichnet. Der Abstand $r$ sei groß genug, sodass $\op H_1$ als Störungsterm behandelt werden kann. Berechnen Sie zuerst die Verschiebung der Energieniveaus $\ket{++}$, $\ket{--}$ in erster Ordnung von $\op H_1$. Zur Berechnung der relevanten Matrixelemente $\bra{\pm\pm}\op H_1\ket{\pm\pm}$ müssen nur die Terme in $\op H_1$ betrachtet werden, die proportional zu $S_z^{(1)}S_z^{(2)}$ sind. Warum? Es ist nützlich, die Frequenz
        \begin{align*}
            \Omega=\frac{\mu_0\gamma^2\hbar}{4\pi r^3}
        \end{align*}
        zu definieren.

        \begin{calculationbox}
            Wir setzen $\vec M^{(i)}=\gamma\vec S^{(i)}$ ein. Dann wird
            \begin{align*}
                \op H_1
                =
                \frac{\mu_0\gamma^2}{4\pi r^3}
                \left[
                    \vec S^{(1)}\cdot\vec S^{(2)}
                    -3(\vec S^{(1)}\cdot\vec n)(\vec S^{(2)}\cdot\vec n)
                \right].
            \end{align*}
            Mit
            \begin{align*}
                \Omega=\frac{\mu_0\gamma^2\hbar}{4\pi r^3}
            \end{align*}
            gilt also
            \begin{align*}
                \frac{\mu_0\gamma^2}{4\pi r^3}=\frac{\Omega}{\hbar}.
            \end{align*}

            \begin{calculationbox}[title=Warum nur $S_z^{(1)}S_z^{(2)}$ beiträgt]
                Die Zustände $\ket{++}$ und $\ket{--}$ sind Eigenzustände von $S_z^{(1)}$ und $S_z^{(2)}$. Ein diagonales Matrixelement wie $\bra{++}\op H_1\ket{++}$ ist nur dann ungleich null, wenn der Operator den Zustand nicht in einen orthogonalen Zustand überführt. Die Operatoren $S_x$ und $S_y$ enthalten Leiteroperatoren,
                \begin{align*}
                    S_x=\frac12(S_++S_-),
                    \qquad
                    S_y=\frac{1}{2i}(S_+-S_-),
                \end{align*}
                und ändern Spinrichtungen. Daher liefern sie bei diesen diagonalen Matrixelementen keinen Beitrag. Es bleibt nur der Anteil proportional zu $S_z^{(1)}S_z^{(2)}$.
            \end{calculationbox}

            \begin{calculationbox}[title=Relevanter Anteil des Störoperators]
                Aus
                \begin{align*}
                    \vec S^{(1)}\cdot\vec S^{(2)}
                    &=
                    S_x^{(1)}S_x^{(2)}+S_y^{(1)}S_y^{(2)}+S_z^{(1)}S_z^{(2)},\\
                    (\vec S^{(1)}\cdot\vec n)(\vec S^{(2)}\cdot\vec n)
                    &\supset
                    \cos^2\theta\,S_z^{(1)}S_z^{(2)}
                \end{align*}
                folgt
                \begin{align*}
                    \op H_1
                    \supset
                    \frac{\Omega}{\hbar}
                    (1-3\cos^2\theta)S_z^{(1)}S_z^{(2)}.
                \end{align*}
            \end{calculationbox}

            \begin{calculationbox}[title=Energieverschiebungen von $\ket{++}$ und $\ket{--}$]
                Für beide Zustände ist
                \begin{align*}
                    \bra{++}S_z^{(1)}S_z^{(2)}\ket{++}
                    &=
                    \frac{\hbar}{2}\frac{\hbar}{2}
                    =
                    \frac{\hbar^2}{4},\\
                    \bra{--}S_z^{(1)}S_z^{(2)}\ket{--}
                    &=
                    \left(-\frac{\hbar}{2}\right)
                    \left(-\frac{\hbar}{2}\right)
                    =
                    \frac{\hbar^2}{4}.
                \end{align*}
                Also folgt
                \begin{align*}
                    \Delta E_{++}^{(1)}
                    =
                    \Delta E_{--}^{(1)}
                    =
                    \frac{\Omega}{\hbar}(1-3\cos^2\theta)\frac{\hbar^2}{4}.
                \end{align*}

                \begin{solutionbox}[title=Verschiebung der äußeren Niveaus]
                    \begin{align*}
                        \Delta E_{++}^{(1)}
                        =
                        \Delta E_{--}^{(1)}
                        =
                        \frac{\hbar\Omega}{4}(1-3\cos^2\theta).
                    \end{align*}
                \end{solutionbox}
            \end{calculationbox}
        \end{calculationbox}

        \item Berechnen Sie nun die Verschiebung der Energieniveaus $\ket{+-}$, $\ket{-+}$ in erster Ordnung von $\op H_1$ und bestimmen Sie die zugehörigen Spinwellenfunktionen. Zur Berechnung der relevanten Matrixelemente $\bra{\pm\mp}\op H_1\ket{\pm\mp}$ bzw. $\bra{\pm\mp}\op H_1\ket{\mp\pm}$ müssen nur die Terme in $\op H_1$ betrachtet werden, die proportional zu $S_z^{(1)}S_z^{(2)}$ bzw. $S_x^{(1)}S_x^{(2)}$, $S_y^{(1)}S_y^{(2)}$ sind. Warum?

        \begin{calculationbox}
            Die Zustände $\ket{+-}$ und $\ket{-+}$ sind bezüglich $\op H_0$ entartet. Daher müssen wir die Störung im entarteten Unterraum
            \begin{align*}
                \mathcal H_0=\operatorname{span}\{\ket{+-},\ket{-+}\}
            \end{align*}
            diagonalisisieren.

            \begin{calculationbox}[title=Warum genau diese Terme beitragen]
                Für diagonale Matrixelemente $\bra{+-}\op H_1\ket{+-}$ und $\bra{-+}\op H_1\ket{-+}$ muss der Zustand unverändert bleiben. Dafür trägt nur $S_z^{(1)}S_z^{(2)}$ bei. Für außerdiagonale Matrixelemente wie $\bra{+-}\op H_1\ket{-+}$ müssen beide Spins gleichzeitig umgedreht werden. Genau diese Kopplung liefern $S_x^{(1)}S_x^{(2)}$ und $S_y^{(1)}S_y^{(2)}$. Terme mit nur einem $S_z$ und einem transversalen Spinoperator ändern höchstens einen Spin und geben wegen Orthogonalität keinen Beitrag.
            \end{calculationbox}

            \begin{calculationbox}[title=Diagonale Matrixelemente]
                Der relevante diagonale Anteil ist
                \begin{align*}
                    \frac{\Omega}{\hbar}(1-3\cos^2\theta)S_z^{(1)}S_z^{(2)}.
                \end{align*}
                Für $\ket{+-}$ gilt
                \begin{align*}
                    \bra{+-}S_z^{(1)}S_z^{(2)}\ket{+-}
                    =
                    \frac{\hbar}{2}\left(-\frac{\hbar}{2}\right)
                    =
                    -\frac{\hbar^2}{4}.
                \end{align*}
                Dasselbe gilt für $\ket{-+}$. Wir definieren zur Abkürzung
                \begin{align*}
                    B:=\frac{\hbar\Omega}{4}(3\cos^2\theta-1).
                \end{align*}

                \begin{solutionbox}[title=Diagonale Matrixelemente]
                    \begin{align*}
                        \bra{+-}\op H_1\ket{+-}
                        =
                        \bra{-+}\op H_1\ket{-+}
                        =
                        B,
                        \qquad
                        B=\frac{\hbar\Omega}{4}(3\cos^2\theta-1).
                    \end{align*}
                \end{solutionbox}
            \end{calculationbox}

            \begin{calculationbox}[title=Außerdiagonale Matrixelemente]
                Die Koeffizienten vor den relevanten Operatoren sind
                \begin{align*}
                    S_x^{(1)}S_x^{(2)} &: \quad 1-3\sin^2\theta\cos^2\varphi,\\
                    S_y^{(1)}S_y^{(2)} &: \quad 1-3\sin^2\theta\sin^2\varphi.
                \end{align*}
                Außerdem gilt
                \begin{align*}
                    \bra{+-}S_x^{(1)}S_x^{(2)}\ket{-+}
                    =
                    \frac{\hbar^2}{4},
                    \qquad
                    \bra{+-}S_y^{(1)}S_y^{(2)}\ket{-+}
                    =
                    \frac{\hbar^2}{4}.
                \end{align*}
                Deshalb
                \begin{align*}
                    \bra{+-}\op H_1\ket{-+}
                    &=
                    \frac{\Omega}{\hbar}\frac{\hbar^2}{4}
                    \left[
                        1-3\sin^2\theta\cos^2\varphi
                        +1-3\sin^2\theta\sin^2\varphi
                    \right]\\
                    &=
                    \frac{\hbar\Omega}{4}(2-3\sin^2\theta)\\
                    &=
                    \frac{\hbar\Omega}{4}(3\cos^2\theta-1)
                    =B.
                \end{align*}

                \begin{solutionbox}[title=Außerdiagonale Matrixelemente]
                    \begin{align*}
                        \bra{+-}\op H_1\ket{-+}
                        =
                        \bra{-+}\op H_1\ket{+-}
                        =
                        B.
                    \end{align*}
                \end{solutionbox}
            \end{calculationbox}

            \begin{calculationbox}[title=Störmatrix und Diagonalisierung]
                Im entarteten Unterraum lautet die effektive Störmatrix
                \begin{align*}
                    H_1^{\mathrm{eff}}
                    =
                    \begin{pmatrix}
                        B&B\\
                        B&B
                    \end{pmatrix}.
                \end{align*}
                Die Matrix hat den symmetrischen Eigenvektor mit Eigenwert $2B$ und den antisymmetrischen Eigenvektor mit Eigenwert $0$.

                \begin{solutionbox}[title=Aufgespaltene Zustände und Verschiebungen]
                    \begin{align*}
                        \ket{\psi_s}
                        =
                        \frac{1}{\sqrt2}(\ket{+-}+\ket{-+}),
                        \qquad
                        \Delta E_s^{(1)}
                        =
                        2B
                        =
                        \frac{\hbar\Omega}{2}(3\cos^2\theta-1),\\[0.6em]
                        \ket{\psi_a}
                        =
                        \frac{1}{\sqrt2}(\ket{+-}-\ket{-+}),
                        \qquad
                        \Delta E_a^{(1)}=0.
                    \end{align*}
                \end{solutionbox}
            \end{calculationbox}
        \end{calculationbox}

        \item Skizzieren Sie, wie sich das Energiespektrum durch die Wechselwirkung verschiebt.

        \begin{calculationbox}
            Ohne Wechselwirkung gibt es die Niveaus $-\hbar\omega_0$, $0$ und $+\hbar\omega_0$, wobei das mittlere Niveau zweifach entartet ist. Aus den vorherigen Teilaufgaben folgt für die Verschiebung der äußeren Niveaus
            \begin{align*}
                A:=\frac{\hbar\Omega}{4}(1-3\cos^2\theta).
            \end{align*}
            Das mittlere Niveau spaltet in einen symmetrischen Zustand mit Verschiebung $-2A$ und einen antisymmetrischen Zustand mit Verschiebung $0$ auf.

            \begin{solutionbox}[title=Gestörtes Spektrum bis erster Ordnung]
                \begin{align*}
                    E_{++}
                    &=
                    -\hbar\omega_0
                    +
                    \frac{\hbar\Omega}{4}(1-3\cos^2\theta),\\
                    E_{--}
                    &=
                    +\hbar\omega_0
                    +
                    \frac{\hbar\Omega}{4}(1-3\cos^2\theta),\\
                    E_s
                    &=
                    \frac{\hbar\Omega}{2}(3\cos^2\theta-1),\\
                    E_a
                    &=
                    0.
                \end{align*}
            \end{solutionbox}

            \begin{calculationbox}[title=Beschreibung der Skizze]
                Für die Skizze zeichnet man links zuerst die ungestörten Energieniveaus:
                \begin{align*}
                    -\hbar\omega_0,
                    \qquad
                    0 \text{ (zweifach entartet)},
                    \qquad
                    +\hbar\omega_0.
                \end{align*}
                Rechts zeichnet man nach Einschalten von $\op H_1$ die verschobenen Niveaus: $\ket{++}$ und $\ket{--}$ verschieben sich beide um $A$. Das vorher entartete Niveau bei $0$ spaltet auf in $E_s=-2A$ und $E_a=0$. Das Vorzeichen der Verschiebung hängt vom Winkel $\theta$ ab.

                \begin{solutionbox}[title=Magischer Winkel]
                    \begin{align*}
                        3\cos^2\theta-1=0
                        \quad\Longleftrightarrow\quad
                        \theta=\arccos\left(\frac{1}{\sqrt3}\right).
                    \end{align*}
                \end{solutionbox}
                Bei diesem Winkel verschwinden die ersten Ordnungsverschiebungen.
            \end{calculationbox}
        \end{calculationbox}
    \end{enumerate}

\end{document}
